\documentclass[10pt,a4paper]{amsart}
\usepackage[margin=19mm]{geometry}
\usepackage{amsmath,amssymb,mathtools}
\usepackage[scr=rsfs,cal=euler]{mathalfa}
\usepackage{xfrac}
\usepackage[unicode,hidelinks,bookmarks=false]{hyperref}
\newcommand{\ie}{i.e.\ }
\newcommand{\cf}{cf.\ }
\newcommand{\resp}{respectively\ }
\newcommand{\Subsection}{\subsection}
\providecommand{\nobreakdash}{\nobreak\hbox{-}}
\setlength{\emergencystretch}{2em}
\multlinegap0pt
\usepackage{bm}
\usepackage{bbm}
\usepackage{dsfont}
\usepackage{xspace}
\usepackage{cancel}
\usepackage{mathtools}
\usepackage{amsthm}
\makeatletter
\@ifundefined{theorem}{\newtheorem{theorem}{Theorem}}{}
\@ifundefined{remark}{\newtheorem{remark}[theorem]{Remark}}{}
\makeatother
\title[An Invariance Principle of 1D KPZ with Robin Boundary Condit]{An Invariance Principle of 1D KPZ with Robin Boundary Conditions}
\author{Yiming Tang}
\date{Source: arXiv:2404.19215v2 (CC BY 4.0)}
\begin{document}
\maketitle

\noindent\emph{Source and licence.} An Invariance Principle of 1D KPZ with Robin Boundary Conditions. Version arXiv:2404.19215v2 (CC BY 4.0), distributed under the Creative Commons Attribution 4.0 International licence, as recorded in the arXiv licence metadata for this exact version and shown on the abstract page https://arxiv.org/abs/2404.19215v2. Full rights notice: licenses/arxiv-2404.19215-CC-BY-4.0.txt.\par\smallskip

\noindent\textit{Editorial scope.} Counted unit: the Remark whose source label is \texttt{None} in the article (source0.tex lines 240--242, source label \texttt{None}), delivered together with the article's own complete original proof of it (source0.tex lines 244--294, one \texttt{proof} environment of 51 lines). No mathematics is written, repaired or re-derived: statement and proof are the article's own text line for line. Required context retained uncounted: 5.1 (lines 234--237). Every definition, display and label of those lines is retained unchanged. Omitted from this package and not redistributed: 1--233, 238--239, 243--243, 295--617. Nothing inside a retained excerpt is omitted or altered: every retained body is delivered exactly as the article's lines are written, so no omission receipt of any kind accompanies this package. Citations: the retained excerpts carry no citation command at all, so no bibliography entry of the article is transcribed and no reference is redistributed. Reference adaptation: none was needed and none was made; every reference inside the delivered unit resolves inside it, so no unresolved reference or citation is produced. Printed slips kept literal: no printed word, symbol, hypothesis or number of the counted statement or of its proof was altered. Layout and class: the article's own class is \texttt{article} with packages including geometry, inputenc, fontenc, amsmath, amsthm, amssymb, amsfonts, yfonts, mathrsfs, float, listings, longtable, tikz-cd, breqn; the wrapper is the lane's pinned \texttt{amsart} editorial wrapper at 10pt on A4, which restates the theorem-like environments the retained text needs and adds the article's own macro definitions that the retained text needs (none), with extra wrapper packages (bm, bbm, dsfont, xspace, cancel, mathtools). The delivered numbering is the unit's own. Assets: no retained span contains a figure environment, a graphics-inclusion command, a PDF-page inclusion, a table or a TikZ picture; no image file is copied into this package, the package declares no asset dependency and no image file is redistributed with it. Licence: version arXiv:2404.19215v2 of this article is distributed under the Creative Commons Attribution 4.0 International licence, as recorded in the arXiv licence metadata for this exact version and shown on the abstract page https://arxiv.org/abs/2404.19215v2. A full rights notice is delivered at licenses/arxiv-2404.19215-CC-BY-4.0.txt. Characters: every retained excerpt is pure ASCII, so the delivered engine, XeLaTeX with its default Latin Modern text font, reports no missing character.
\begin{theorem} \label{5.1}
    Fix any $a, b >0$, there exists $C_1, C_2 >0$ such that for all sufficiently large $N$,  $(x,t) \in [0,a\sqrt{N}] \times [0,bN]$, 
    $$C_1 \leq \frac{1}{2^t} \sum_{q \in RW(x,t)} \gamma^{d(q)} \leq C_2.$$
\end{theorem}

\begin{remark}
Note that $\mathbb{E}[{Z(x,t)}]=\frac{1}{2^t}\sum_{q \in RW(x, t)} \Lambda(|q(0)|) \gamma^{d(q)}$. Under our assumption on the initial condition that $C^{-1} \leq \Lambda_N(x) \leq C$ for a  constant $C$ uniform over $N$, it thus suffices to bounds the exponential moment of the local time with a flat initial condition. Combining with Theorem \ref{5.1}, we will obtain  lower and upper bounds uniformly on $[0,a\sqrt{N}] \times [0,bN]$ for $\mathbb{E}[{Z(x,t)}]=\frac{1}{2^t}\sum_{q \in RW(x, t)} \Lambda(|q(0)|) \gamma^{d(q)}$.
\end{remark}

\begin{proof}

Define $TRW(x,t)$ to be the set of simple symmetric random walks on $[0, t]$ starting from $x$, which corresponds to the time reversal of $RW(x,t)$. Let $\tilde{d}(q)$ be the number of visits to 0 of a simple symmetric random walk $q \in TRW(x, t)$ in time $[0, t-1]$. Then the exponential moment of the local time can also be expressed as $\frac{1}{2^{t}}\sum_{q \in TRW(x,t)} \gamma^{\tilde{d}(q)}.$


By Jensen's inequality, as $f(w)=\gamma^w$ is convex, $\frac{1}{2^{t}}\sum_{q \in TRW(x,t)} \gamma^{\tilde{d}(q)} =\mathbb{E}[\gamma^{\tilde{d}}] \geq \gamma^{\mathbb{E}[\tilde{d}]}.$

It is a standard result that $\mathbb{E}[\tilde{d}] \approx \sqrt{t}$, which the $\approx$ sign means there exist non-zero universal constants bounding the ratio of the two sides. So there exist a constant $c>0$ such that $\frac{1}{2^{t}}\sum_{q \in TRW(x,t)} \tilde{d}(q) \leq c\sqrt{t} \leq c\sqrt{bN}$. For $\gamma <1$, the exponential moment of the local time is lower bounded by
$(1-\frac{A}{\sqrt{N}})^{c\sqrt{bN}},$
which converges to $e^{-Ac\sqrt{b}} >0$. For $\gamma\geq 1$, we obtain a trivial lower bound of 1.

To prove the other direction, for any $s \in [1,t]$,
 $|q(s)|-|q(s-1)| = \text{sgn}(q(s-1))(q(s)-q(s-1))+\delta_0(q(s-1))$.

Summing over $s=1$ to $t$, we get
\begin{equation} 
    |q(t)|-|q(0)|=\sum_{s=1}^t \text{sgn}(q(s-1))(q(s)-q(s-1)) + \sum_{s=1}^t \delta_0(q(s-1)).
\end{equation}
Given $|q(0)|=x$, we have
\begin{equation} \label{mean}
    \frac{1}{2^t}\sum_{q \in TRW(x,t)} \gamma^{\tilde{d}(q)}=\frac{1}{2^{t}}\sum_{q \in TRW(x,t)} \gamma^{|q(t)|-x-\sum_{s=1}^t \text{sgn}(q(s-1))(q(s)-q(s-1))}.
\end{equation}

Note that $\gamma^{-x}$ is bounded uniformly over $[0, a\sqrt{N}]$ both above and below by positive constants independent of $N$ as $x=O(\sqrt{N})$ and $\gamma =1-O(N^{-1/2})$. By Cauchy-Schwartz, the rest of (\ref{mean}) is bounded above by 
\begin{equation}
    \Big(\frac{1}{2^{t}}\sum_{q \in TRW(x,t)} \gamma^{2|q(t)|}\Big)^{1/2}\Big(\frac{1}{2^{t}}\sum_{q \in TRW(x,t)} \gamma^{-2\sum_{s=1}^t \text{sgn}(q(s-1))(q(s)-q(s-1))}\Big)^{1/2}.
\end{equation}

For any $x \in [0, a\sqrt{N}]$, and $t \in [0, bN]$, we have
\begin{equation}
        \frac{1}{2^{t}}\sum_{q \in TRW(x,t)} \gamma^{2|q(t)|} =\frac{1}{2^t} \sum_{k=0}^t {{t} \choose {k}} \gamma^{2|x+2k-t|} \leq \frac{1}{2^t} \sum_{k=0}^t {{t} \choose {k}} \left(\gamma^{2(x+2k-t)}+\gamma^{2(-x-2k+t)}\right),
\end{equation}
which equals to $2^{-t}\left(\gamma^{-2}+\gamma^2\right)^t\left(\gamma^{2x}+\gamma^{-2x}\right)$ and therefore it is bounded above uniformly over $(x, t) \in [0, a\sqrt{N}] \times [0, bN]$.

For $\mathbb{E}[\gamma^{-2\sum_{s=1}^t \text{sgn}(q(s-1))(q(s)-q(s-1))}]$, we adopt a standard martingale argument. Take $q \in TRW(x,t)$ and define $\mathcal{F}_s:=\sigma(\{q(i): 1 \leq i \leq s\})$. Then the expectation equals to
\begin{equation}
    \mathbb{E}[\gamma^{-2\sum_{s=1}^{t-1} \text{sgn}(q(s-1))(q(s)-q(s-1))} \mathbb{E}[\gamma^{-2\text{sgn}(q(t-1))(q(t)-q(t-1))}|F_{t-1}]],
\end{equation}
by the tower property.

Note that $\mathbb{E}[\gamma^{-2\text{sgn}(q(t-1))(q(t)-q(t-1))}|F_{t-1}]=\frac{\gamma^2+\gamma^{-2}}{2}(1-\delta_0(q(t-1)))+\delta_0(q(t-1)) \leq \frac{1}{2}\left(\gamma^2+\gamma^{-2}\right)$, which is $1+O(N^{-1})$.

Inductively perform the step above, as $t \leq bN$ and $x \in [0, a\sqrt{N}]$, we would eventually get 
\begin{equation}
    \mathbb{E}\left[\gamma^{-2\sum_{s=1}^t \text{sgn}(q(s-1))(q(s)-q(s-1))}\right] =(1+O(N^{-1}))^t \lesssim 1.
\end{equation}
    

Combining the inequalities above, we have shown the existence of a constant upper bound for the exponential moment of the local time.

\end{proof}

\end{document}
